% Folder 136, batch 03 — archivists' pages 41–60.
% Pass: Opus 5.5 (claude-opus-5-5), 2026-10-03 — first pass, unchecked against the pages by a human.
%
% Pages 41, 43, 45, 46, 48, 50, 52, 54, 56, 58 and 60 are not transcribed:
% typed versos by other hands (administrative minutes, typescripts) or blank.
% Pages 47, 49 and 51 are in very faint pencil; their prose is largely lost.
\documentclass[11pt,a4paper]{article}
\input{../preamble/grothendieck}

\folder{136}
\batch{3}
\pages{41}{60}
\foldertitle{Complexe de De Rham à puissance divisée [conférence de 1976 à l’IHÉS] : notes manuscrites (s.d.).}
\dating{[à partir de 1975-1976]}
\watermark{Édition de démonstration}

\begin{document}

\section{Puissances divisées d'une variable : les $T_r = T^{(p^r)}$}

\page{42}
\note{la page continue un calcul commencé avant ce lot : $S_k = k\{T\}$, $\tau_{p^r}S_k$, les homomorphismes de degré $-d$ ne sont pas définis ici.}

\marginal{anneau de base $k$ une $\mathbb{Z}_p$-algèbre}

$T_r$ \quad $(= T^{(p^r)})$
\[
  T_r^{\,p} = p!\,\bigl((T^{(p^r)})^{(p)}\bigr) = p!\,\gamma_r\, T^{(p^{r+1})} = p\,\nu'_r\, T_{r+1},
  \qquad \gamma_r = \frac{p^{r+1}!}{(p^r!)^p\,p!} \in \mathbb{Z}_p^{*},
\]
\[
  \nu'_r = (p-1)!\,\frac{p^{r+1}!}{(p^r!)^p\,p!} \in \mathbb{Z}_p^{*} .
\]
Encadré :
\[
  \boxed{T_r^{\,p} = p\,\nu'_r\, T_{r+1}}, \quad \nu'_r \in \mathbb{Z}_p^{*}
  \qquad\qquad \boxed{S_k = k\{T\}}
\]

$\tau_{p^r} S_k$ a comme base les $T^{\underline{i}}$, $\underline{i} = (i_0, i_1, \ldots, i_r, i_{r+1}, \ldots)$,
où $\sum i_\alpha p^\alpha \geq p^r$ i.e. $\exists\, \alpha \geq r$ avec $i_\alpha \neq 0$.

Donc $\tau_{p^r} S$ est engendré par $T_r, T_{r+1}, \ldots, T_{r+\ell}, \ldots$ ($\ell \geq 0$),
notés $t_r, t_{r+1}, \ldots$ ; les relations sont
\[
  \boxed{T_s^{\,p-1}\, t_s = p\,\nu'_s\, t_{s+1} \quad \text{si } s \geq r .}
\]

Donc les \uncertain{homomorphismes} de degré $d < p^r$, $d = d_0 + d_1 p + \cdots + d_{r-1}p^{r-1}$,
\[
  u : \tau_{p^r} S_k \longrightarrow M \otimes_k S_k
\]
\note{devant $M$ une lettre biffée.}
est donné par les
\[
  \boxed{u(t_s) = \eta_s \in S^{(p^s - d)} \otimes_k M}
\]
\note{un premier second membre, $p^{s}-d$ \ldots, est biffé ; à gauche, un repère biffé.}
avec les relations
\[
  (R_s) \quad \boxed{T_s^{\,p-1}\,\eta_s = p\,\nu'_s\,\eta_{s+1}}
\]

Soit $\nu = v_p(d)$, de sorte que
\[
  d = d_\nu p^\nu + \cdots + d_{r-1}p^{r-1} = p^\nu\bigl(d_\nu + d_{\nu+1}p + \cdots + d_{r-1}p^{r-1-\nu}\bigr),
  \qquad d_\nu \neq 0 \text{ i.e. } d_\nu \geq 1 .
\]
\begin{align*}
  p^s - d &= p^\nu\Bigl(p^{s-\nu} - \frac{d}{p^\nu}\Bigr)
           = p^\nu\Bigl(p^{s-\nu} - 1 - \underbrace{\Bigl(\frac{d}{p^\nu} - 1\Bigr)}_{(d_\nu - 1) + d_{\nu+1}p + \cdots + d_{r-1}p^{r-1-\nu}}\Bigr) \\
          &= p^\nu\bigl(\delta_\nu + \delta_{\nu+1}p + \cdots + \delta_{r-1}p^{r-\nu-1} \\
          &\qquad + \underbrace{(p-1)p^{r-\nu} + (p-1)p^{r-\nu+1} + \cdots + (p-1)p^{s-\nu-1}}_{\text{si } s > r}\bigr),
\end{align*}
avec
\[
  \delta_\nu = (p-1) - (d_\nu - 1), \quad \delta_{\nu+1} = (p-1) - d_{\nu+1}, \quad \ldots, \quad \delta_{r-1} = (p-1) - d_{r-1} .
\]

Donc $S^{(p^s-d)}$ a comme base
\[
  T_\nu^{\delta_\nu}\, T_{\nu+1}^{\delta_{\nu+1}} \cdots T_{r-1}^{\delta_{r-1}}\,\bigl(T_r \cdots T_{s-1}\bigr)^{p-1}
\]
\uncertain{à facteur près}, $s > r$ \ill{}.
\note{devant le premier facteur, un $T$ biffé.}

\page{44}
On pourra donc poser
\[
  \eta_s = \xi_s\, T_\nu^{\delta_\nu}\, T_{\nu+1}^{\delta_{\nu+1}} \cdots T_{r-1}^{\delta_{r-1}}\,\bigl(T_r \cdots T_{s-1}\bigr)^{p-1}
  \qquad \text{avec } \xi_s \in M .
\]
On aura
\[
  \eta_s T_s^{\,p-1} = \xi_s\, T_\nu^{\delta_\nu} \cdots T_{r-1}^{\delta_{r-1}}\,\bigl(T_r \cdots T_s\bigr)^{p-1} .
\]
La relation $(R_s)$ s'écrit donc, comme
\[
  p\,\nu'_s\,\eta_{s+1} = p\,\nu'_s\,\xi_{s+1}\, T_\nu^{\delta_\nu} \cdots T_{r-1}^{\delta_{r-1}}\,\bigl(T_r \cdots T_s\bigr)^{p-1},
\]
\[
  \boxed{\xi_s = p\,\nu'_s\,\xi_{s+1}}
\]
On aura donc
\[
  \xi_s = p^h\, \nu'_s \cdots \nu'_{s+h-1}\, \xi_{s+h} \qquad \forall\, h \geq 1,
\]
donc
\[
  \xi_s \in \bigcap_h \bigl(p^h M\bigr),
\]
\note{devant $\bigcap$ un signe biffé.}
i.e.\ les $\xi_s$ \uncertain{appartiennent} à la partie divisible de $M$.

\textbf{Prop.} Soient $k$ un anneau commutatif, $S_k = k\{T\}$, $M$, $N$ deux $k$-modules,
\struck{avec} $d > 0$, $N_0^{\natural} = \bigcap_{n \in \mathbb{N}^{*}} nN$ le \struck{partie}
\add{plus grand sous-module} divisible de $N$, $\alpha \geq d$. Alors si
$u \in \mathrm{Hom}^{-d}(\tau_\alpha M_{S_k}, N_{S_k})$
(\struck{\ill} \add{lequel modulo} \uncertain{est indépendant} \uncertain{modulo restrictions} du choix de $\alpha$),
$u$ se factorise par $N_{0\,S_k}$, i.e.
\[
  \mathrm{Hom}^{-d}(\tau_\alpha M_{S_k}, N_{0\,S_k}) \xrightarrow{\ \sim\ } \mathrm{Hom}^{-d}(\tau_\alpha M_{S_k}, N_{0\,S_k}) .
\]
\note{les deux membres sont écrits identiques sur la page ; le contexte demande $N_{S_k}$ au second.}

\emph{Cor.} Si $N_0 = 0$, p.~ex.\ $k$ de torsion, alors $\mathrm{Hom}^{-d}(\tau_\alpha M_{S_k}, N_{S_k}) = 0$.

\section{Cohomologie des ensembles semi-simpliciaux}

\page{47}
\note{page au crayon, très pâle ; seules les formules et quelques mots se lisent.}

\marginal{Cohomologie des ens.\ semi-simpliciaux et résolutions canon.\ \uncertain{semi-simpliciales}}

Soit $k$ un \struck{groupe abélien} \add{anneau}, $M$ un $k$-module,
\[
  C^{*} : \widehat{\Delta}^{\circ} \longrightarrow \ldots
\]
(complexes \ill{} de $k$-modules) un foncteur \ill{} tel que
\begin{itemize}
  \item[a)] $\forall\, X_\bullet \in \mathrm{Ob}\,\widehat{\Delta}$, $H^{*}(X_\bullet, M) \simeq H^{*}(C^{*}(X_\bullet))$
  (\uncertain{isom.\ fonctoriel en $M$} \ill{}) ;
  \item[b)] $\forall\, n \geq 0$, $X_\bullet \mapsto C^{n}(X_\bullet)$ \uncertain{transforme} \ill{} en \ill{}.
\end{itemize}
Alors b) signifie que les $C^{n}$ sont représentables :
\[
  C^{n}(X_\bullet) = \mathrm{Hom}(X_\bullet, C^{n}_\bullet) .
\]
\ill{} les $C^{*}_\bullet$ \uncertain{sont des} \uncertain{résolutions} de $M$, \ill{}
\[
  M_\bullet = \underline{k} \otimes M \ldots
\]
\ill{} $H^{*}(C^{*}(X_\bullet)) \to H^{*}(X_\bullet, M)$ \ill{}

\note{suit une dizaine de lignes illisibles, un petit diagramme $C^{0}_\bullet \to Z^{1}_\bullet$, $X_\bullet \to Z^{1}_\bullet$, et un passage biffé par des diagonales ; puis :}
\ill{} On aura
\[
  H^{0}(X_\bullet, C^{0}_\bullet) \to H^{0}(X_\bullet, Z^{1}_\bullet) \to H^{1}(X_\bullet, M)
  \to H^{1}(X_\bullet, C^{0}_\bullet) \to H^{1}(X_\bullet, Z^{1}_\bullet)
\]
Mais \ill{} \ill{}

\emph{NB} \ill{} $C^{*}_\bullet$ \ill{} $M_\bullet$, les $Z^{n}(C^{*}_\bullet)$ \ill{} $K(M, n)$ …

\page{49}
\note{page au crayon, très pâle.}

\textbf{Ex.~1.} \ill{}
\[
  M \mapsto C^{*}(K_\bullet, M) \overset{\text{déf}}{=} \mathrm{Hom}_{\widehat{\Delta}}(K_\bullet, M)
  = \mathrm{Hom}(\mathbb{Z}^{(K_\bullet)}, M) \simeq \mathrm{Hom}_{\widehat{\Delta}}(K_\bullet, \ldots)
\]
est représentable \ill{}. \ill{} $\widehat{\Delta}$ \ill{}

\ill{} $C^{*}_{\bullet M}$. \ill{}
\[
  C^{*}(K_\bullet, M) \simeq \mathrm{Hom}(K_\bullet, C^{*}_{\bullet M})
\]
\ill{} $k$-modules \ill{} $(K_\bullet)$ \ill{} $M_\bullet$, \ill{} $C^{*}_{\bullet M}$ \ill{} \ill{}
\[
  C^{n}_{\bullet M} = \ldots \simeq k^{\mathrm{Hom}(\Delta_n, \Delta_m)} \otimes_k M
\]

\textbf{Ex.~2.} Considérons le sous-complexe \ill{} des cochaînes \ill{} des $C^{*}_{\bullet M}$
formé des cochaînes \ill{}
\[
  C^{n}_{\bullet M} = \bigcap \mathrm{Ker}\bigl(C^{n}_{\bullet M} \to \ldots\bigr) = \ldots
\]
\note{trois intersections de noyaux, indexées par des conditions illisibles.}

On sait que \ill{} $C^{*}(X_\bullet)$ \ill{}

\page{51}
\note{page au crayon, très pâle ; suite de l'Ex.~2.}

\ill{}
\[
  C^{*!}(X_\bullet) = \mathrm{Hom}(X_\bullet, C^{*!}_{\bullet M})
\]
\ill{} explicitement
\begin{align*}
  (C^{n!}_{\bullet M})_m &\simeq M^{\mathrm{Hom\,mono}(\Delta_n, \Delta_m)} \simeq k^{\mathrm{Hom}(\Delta_n, \Delta_m)} \otimes_k M \\
  &\simeq \Bigl(\bigwedge^{n+1} k^{\Delta_m}\Bigr) \otimes_k M
\end{align*}
\ill{}
\[
  C^{n!}_{\bullet M} \simeq \bigwedge^{n+1} \ldots \otimes_k M
\]
\note{le reste de la page (une quinzaine de lignes, où reviennent $\bigwedge$, $\otimes_k M$ et $C^{n!}_{\bullet M}$) ne se lit pas.}

\section{Intégration des algèbres semi-simpliciales}

\page{53}
\textbf{Intégration de \uncertain{algèbres} semi-simpliciales}

Soit $A_{*}$ un groupe gradué différentiel semi-simplicial, i.e.\ un objet semi-simplicial dans
la catégorie des groupes gradués différentiels (à opérateur de dérivation de degré $+1$).
\struck{\ill{}} Donc $\forall\, n \geq 0$, $A_n = A_n^{*} = \sum_{p \in \mathbb{Z}} A_n^{p}$, avec
opér.\ différentielle $d_n : A_n^{i} \to A_n^{i+1}$, $A_n^{p} \to A_n^{p+1}$. On suppose
\begin{itemize}
  \item[(1)] $\forall\, n$, on a $A_n^{p} = 0$ si $p \notin [0, n]$, et \struck{$k$} $A_n^{*}$ ---
  \add{muni} de l'augmentation $\varepsilon_n : A_0 \to A_n^{*}$
  [\emph{NB} \uncertain{définie} : l'unique $\sigma : \Delta_n \to \Delta_0$ ; $A_0$ est considéré
  comme groupe \uncertain{gradué} de \uncertain{graduation} nulle et d'op.\ diff.\ triviale] est une
  \emph{résolution} de $A_0$.
\end{itemize}

Rappelons qu'il y a sur le système des $A_n$ un op.\ différentiel (de degré $-1$)
\[
  \partial : A_n^{*} \to A_{n-1}^{*}, \qquad
  \partial(\omega) = \sum_{0 \leq i \leq n} \partial_{n,i}^{*}(\omega)(-1)^{i}
\]
(où les $\partial_{n,i}$ ($0 \leq i \leq n$) sont les \uncertain{n} \uncertain{applications} strictement
croissantes $\Delta_{n-1} \to \Delta_n$). On sait que $\partial(\partial\omega) = 0$.

\textbf{Proposition 1.} Il existe un et un seul système d'homomorphismes de groupes
\[
  {}^{*}A_n^{*} \xrightarrow{\ K_n\ } A_0
\]
satisfaisant aux conditions suivantes :
\begin{itemize}
  \item[a)] (normalisation) \quad $K_0 = \mathrm{id}_{A_0}$ ;
  \item[b)] (condition de degré) \quad $K_n(A_n^{p}) = 0$ si $p \neq n$ ;
  \item[c)] (Stokes) \quad $K_n(d\omega) = K_{n-1}(\partial\omega)$ \quad $\forall\, \omega \in A_n^{*}$.
\end{itemize}

\marginal{NB le système des $K_n$ \uncertain{est} \ill{} \uncertain{unique} \ill{} de degré \ill{}}

En vertu de b), la condition c) \uncertain{signifiant} \ill{} \ill{} condition, avec
$\omega \in A_n^{n-1}$. On montre construction et unicité \add{degré $\in \mathbb{N}$ par}
\ill{} récurrence sur $N$, en notant que pour $N = 0$ \uncertain{elle est triviale}. Si $N > 0$,
\struck{\ill{} $\in A$} il revient à démontrer \uncertain{existence} et unicité de $K_N$ \ill{}.
Ceci donnera $K_N(\varpi^{N})$ pour $\varpi^{N} \in A_N^{N}$.
Or $d\varpi^{N} = 0$ \struck{\ill{}}, \uncertain{puisque} $A_N^{N+1} = 0$ \add{par (1)}, donc
\[
  \varpi^{N} = d\omega^{N-1} \quad (1)
\]
[En fait, il suffit ici que $H^{N}(A_N^{*}) = 0$ pour $N \geq 1$] on \uncertain{la} \add{ier}
\emph{définit} par c) par
\[
  K_N(d\omega^{N-1}) = K_{N-1}(\partial\omega) .
\]
\note{« ier » est ajouté sous « on la » ; lecture incertaine.}

\page{55}
il faut simplement voir que celui-ci ne dépend pas du choix de $\omega^{N-1}$, i.e.\ que
\struck{$d\omega^{N-1} = 0$}
\[
  d\omega^{N-1} = 0 \;\Longrightarrow\; K_{N-1}(\partial\omega^{N-1}) = 0 .
\]
Distinguons deux cas :

a) $N \geq 2$ i.e.\ $N - 1 \geq 1$, alors $H^{N-1}(A_N^{*}) = 0$ par (1), donc
$d\omega^{N-1} = 0 \Rightarrow \omega^{N-1} = d\alpha^{N-2}$, et on aura
\[
  K_{N-1}(\partial\omega^{N-1}) = K_{N-1}(\partial d\alpha^{N-2})\ldots
\]
Or \struck{$\partial\alpha^{N-}$} $\partial d\alpha^{N-2} = d\partial\alpha^{N-2}$ donc \ill{}
\[
  K_{N-1}(\partial\omega^{N-1}) = K_{N-1}(d\partial\alpha^{N-2})
  \overset{\text{Stokes}}{=} K_{N-2}\bigl(\underbrace{\partial\partial\alpha^{N-2}}_{0}\bigr) = 0
\]
\note{l'indice de la dernière $K$ est mal formé ; $N-2$ est ce que demande Stokes.}
cqfd

b) $N = 1$ i.e.\ $N - 1 = 0$. Alors en vertu de (1) on a :
$d\omega^{0} = 0$ implique $\omega^{0} = \varepsilon_1(\alpha)$ ($\alpha \in A_0$), donc
\[
  K_{N-1}(\partial\omega^{0}) = K_{N-1}\bigl(\underbrace{\partial\varepsilon_1(\alpha)}_{=0}\bigr) = 0
\]
\struck{or $\partial\varepsilon_1(\alpha)$} o.k.
\note{les exposants de $\omega$ sont surchargés par des taches d'encre ; « 0 » est lu d'après la suite.}

\emph{Remarques.} 1) Dans (1) \add{\uncertain{libellé}} on \uncertain{peut} \uncertain{se passer} de la
condition des degrés, \uncertain{on} n'utilise \ill{}
\[
  \begin{cases}
    H^{N}(A_N^{*}) = 0 & \text{si } N \geq 1 \\[2pt]
    H^{N-1}(A_N^{*}) = \begin{cases} 0 & \text{si } N \geq 2 \\ \varepsilon_1(A_0) & \text{si } N = 1 \end{cases}
  \end{cases}
\]

2) On peut donner une version « abstraite » en termes d'un bicomplexe augmenté avec des conditions
de degré \uncertain{finies} ($d\partial = \partial d$ \ill{}) \struck{\ill{}) dans} une cat.\ abélienne
\uncertain{$\mathcal{A}$} [celle pour \uncertain{être} $(\mathrm{Ab})^{\circ}$ !]

\begin{tikzcd}[column sep=small, row sep=small, nodes={font=\scriptsize}]
 & & & & 0 & \\
 & & & 0 & A_3^{3} \arrow[u] \arrow[l] & \cdots \arrow[l] \\
 & & 0 & A_2^{2} \arrow[u] \arrow[l] & A_3^{2} \arrow[u, "d"] \arrow[l, "\partial"'] & \cdots \arrow[l] \\
 & 0 & A_1^{1} \arrow[u] \arrow[l] & A_2^{1} \arrow[u, "d"] \arrow[l, "\partial"'] & A_3^{1} \arrow[u, "d"] \arrow[l, "\partial"'] & \cdots \arrow[l] \\
0 & A_0 \arrow[l] & A_1^{0} \arrow[u, "d"] \arrow[l, "\partial"'] & A_2^{0} \arrow[u, "d"] \arrow[l, "\partial"'] & A_3^{0} \arrow[u, "d"] \arrow[l, "\partial"'] & \cdots \arrow[l] \\
0 \arrow[r] & A_0 \arrow[u, "\varepsilon_0"] & A_0 \arrow[u, "\varepsilon_1"] \arrow[l, "0"'] & A_0 \arrow[u, "\varepsilon_2"] \arrow[l, "\mathrm{id}"'] & A_0 \arrow[u, "\varepsilon_3"] \arrow[l, "0"'] & \cdots \arrow[l, "\mathrm{id}"']
\end{tikzcd}

\note{à côté de $\varepsilon_0$ est écrit « id » ; les flèches $\partial$ des lignes supérieures ne portent pas toutes leur étiquette sur la page.}

conditions :
\struck{$A_N^{*} = 0$ si … ; $H^{N}(A_N^{*}) = 0$ si $N \geq 1$}
\[
  \begin{cases}
    A_N^{i} = 0 & \text{si } i \notin [0, N] \\[2pt]
    H^{N}(A_N^{*}) = 0 & \text{si } N \geq 1 \\[2pt]
    H^{N-1}(A_N^{*}) \begin{cases} = 0 & \text{si } N \geq 2 \\ = \mathrm{Ker}\,\delta & \text{si } N = 1 \end{cases}
  \end{cases}
\]

3) À cause de \uncertain{l'unicité}, on \uncertain{voit} grâce à la \uncertain{version} « intégrale »
\ill{} \uncertain{on voit} que les \ill{} « \uncertain{intégrales} »
$A_n^{n} \to A_0^{0} = A_0$ commutent \ill{} \ill{} les \ill{}
des $A_*^{*}$, \ill{} \ill{} \ill{}.

\section{Topos induit par un ensemble semi-simplicial}

\page{57}
Soit $(SS)$ la catégorie des ensembles semi-simpliciaux
\[
  (SS) = \widehat{\Delta} = \underline{\mathrm{Hom}}(\Delta^{\circ}, (\mathrm{Ens})) .
\]
C'est un topos, et si $X \in \mathrm{Ob}(SS)$ est un ens.\ semi-simplicial, i.e.\ un objet de ce
topos, i.e.\ \uncertain{induit}
\[
  (SS)/X \simeq \widehat{\Delta/X} \simeq \underline{\mathrm{Hom}}\bigl((\Delta/X)^{\circ}, \mathrm{Ens}\bigr) .
\]
Pour nous, l'étude cohomologique de $X$ est celle de ce topos induit. Les faisceaux
\add{(d'ens.)} \struck{\ill{}} \uncertain{sur} « $X$ » sont les \struck{\ill{}} faisceaux
semi-simpliciaux sur $F \to X$ sur $X$, ou ce qui revient au même les foncteurs
$(\Delta/X)^{\circ} \to (\mathrm{Ens})$ ; les faisceaux abéliens sur $X$ sont les objets groupes
abéliens $F \to X$ sur $X$, ou ce qui revient au même les foncteurs
$(\Delta/X)^{\circ} \to (\mathrm{Ab})$, etc. Il \uncertain{en est de même} \struck{f} \uncertain{pour}
faisceaux \uncertain{d'algèbres} graduées différentielles etc. Si $A_\bullet$ est un ens.\ s.s., on peut
regarder
\[
  A \times X \longrightarrow X
\]
comme un faisceau sur $X$ (c'est aussi le faisceau induit par $A_\bullet$ sur le topos induit
$(SS/X)$) ; c'est aussi le faisceau correspondant
\[
  (\Delta/X)^{\circ} \longrightarrow (\Delta)^{\circ} \xrightarrow{\ A\ } (\mathrm{Ens}) .
\]

\page{59}
Itou pour $A_\bullet$ muni d'une structure \uncertain{supplémentaire}, p.~ex.\ de groupe abélien,
éventuellement d'alg.\ graduée. Dans ce dernier cas, si les degrés sont bornés inférieurement
\ill{}, on peut donc parler de l'hypercohomologie
\[
  \mathbb{R}\Gamma(X, A_\bullet^{*}) \overset{\text{déf}}{=} \mathbb{R}^{*}\Gamma_X(\underline{A_X})
\]
\uncertain{complexe induit} par $A_\bullet^{*}$ sur $X$, ou $A^{*} \times X / X$.

Soit $A : (\Delta/X)^{\circ} \to (\mathrm{Ab})$ \uncertain{un groupe sur $X$}, considérons le complexe
des cochaînes sur $X$ à coeff.\ dans $A$ :
\[
  C^{n}(X, A) \overset{\text{déf}}{=} \prod_{\sigma \in \mathrm{Ob}\,\Delta_n/X} A(\sigma),
  \qquad \mathrm{Ob}\,\Delta_n/X = X_n .
\]
Les $C^{n}$ forment un groupe cosimplicial.
\note{la page s'arrête ici ; la construction se poursuit au-delà de ce lot.}

\end{document}
